\honest{Resist the Happy Death Spiral}{Eliezer Yudkowsky}{2007}

Once there was a man sure that his Great Idea would unravel the universe, overturn the Establishment, feed the hungry and heal the sick. He was Francis Bacon, the idea was the scientific method, and he was ``the only crackpot in all history'' to claim that much and turn out ``completely right.'' \nb{Right about the benefits. Bacon also wrote, in the \textsc{Novum Organum} (1620), that if anyone could answer our questions of nature, ``the invention of all causes and sciences would be the labor of but a few years.'' Four centuries later that work is not done. Judged in parts, as the post later advises, his Great Idea held up in some and not in others.} So refusing to admire anything that much cannot be the cure: some ideas really are that good.

How, then, to resist a happy death spiral about science? Not with the standard complaints, such as ``Science gave us air conditioning, but it also made the hydrogen bomb.'' Their originators, I say, ``Probably'' disliked what science said about their pet beliefs. I name none. \nb{Some of the strongest early statements against the hydrogen bomb came from scientists. In 1949 the physicists Enrico Fermi and I.~I. Rabi, advising the US Atomic Energy Commission, called it ``necessarily an evil thing considered in any light.''} Not with a selective search for negatives, which is rationalization. Nor by confining science to a narrow box, since science bears on nearly everything. Martine Rothblatt founded United Therapeutics to seek a cure for her daughter's pulmonary hypertension: ``Successfully, I might add.'' \nb{The company's drugs treat the disease. The US National Heart, Lung, and Blood Institute says: ``There is usually no cure for pulmonary hypertension.''}

What is a false nice claim about science? My first example, ``in my humble opinion,'' is that scientists need not take ethical responsibility because science will turn out well anyway. That point, I admit, is ``not beyond dispute,'' so I switch to simpler ones: publishing enough papers cures a cancer patient; a rule of $p < 0.05$ cures sociopaths. \nb{No admirer of science makes these claims. The method below is shown only on them, although the post itself says the danger lies in ``unsettled'' claims, like the one it set aside.}

The method: know specifically how science works, and ask whether the proposed chain of cause and effect would work. The halo effect ``will probably always influence us a little,'' but you can treat each added nice claim as a burdensome detail and keep the halos subcritical. Do not take happiness from claims that ``can't be disproven.''

Stuart Armstrong, a reader, advises splitting the Great Thingy into parts and judging each independently. That, I say, is like keeping subcritical masses of plutonium apart: ``Three Great Ideas are far less likely to drive you mad than one Great Idea.'' \nb{No evidence is given. The halo effect, as the previous post defined it, acts between separate traits of one person, which is what splitting produces. Judging the parts independently is the hard part, and splitting does not supply it.}

I end with a list: five things that work and three that do not. \nb{The post reports no test of any of them.}
